\section{Related Work}

\paragraph{Exact symmetries and symmetric solutions.}
Symmetry has been part of equilibrium analysis since the classical existence theory for symmetric games \cite{nash1951noncooperative}.  Recent work gives a general computational treatment of normal-form relabelings that may exchange players, actions, or both.  Tewolde et al. connect the discovery of such symmetries to graph automorphism and show that computing a Nash equilibrium constrained to respect a supplied group remains PPAD-complete in general-sum games and CLS-complete in team games, while important zero-sum and high-symmetry cases are tractable \cite{tewolde2025game}.  Symmetric common-payoff bimatrix games can likewise retain substantial equilibrium-computation complexity \cite{ghosh2026complexity}.  These works require exact payoff preservation.  We instead take a semantically proposed group as input, quantify the smallest strategic distortion needed to make it exact, and attach an explicit incentive certificate to the resulting solution.

\paragraph{Strategic equivalence and distances between games.}
The distinction between payoff levels and unilateral incentives is central to game decomposition.  Candogan et al. decompose a finite game into potential, harmonic, and nonstrategic components and show that the nonstrategic component does not affect equilibrium behavior \cite{candogan2011flows}; Hwang and Rey-Bellet similarly study strategic equivalence to zero-sum and potential games \cite{hwang2020strategic}.  The maximum pairwise difference metric was introduced to control how equilibria and learning dynamics transfer from a nearby potential game \cite{candogan2013near}.  Equilibrium-invariant metrics have also been used to organize the geometry of $2\times2$ games \cite{marris2023equilibrium}.  Our metric is the same incentive-level seminorm, but the target set is new: the fixed subspace of an arbitrary player--action relabeling group.  This produces group-specific projection algorithms, critical-cycle obstruction witnesses, and orbit-reduced equilibrium programs.

\paragraph{Approximate symmetry and abstraction.}
Partial and approximate permutation symmetry has been exploited in multi-agent reinforcement learning to trade statistical efficiency against model mismatch \cite{yu2024partial,yardim2025approximately}.  Related abstraction guarantees aggregate approximately equivalent states or information sets in MDPs, Markov games, and extensive-form games \cite{ravindran2004mdp,kroer2014extensive,ishibashi2025approximate}; lossless game abstraction can also be built from exact isomorphisms \cite{gilpin2007lossless}.  Those results act on states, histories, or dynamic models.  Our abstraction acts directly on players, actions, and payoff coordinates of an explicit strategic-form game.  The error is measured by the worst distortion of a unilateral deviation, which is exactly the quantity required for equilibrium transfer.

\paragraph{Correlated and coarse-correlated equilibrium.}
Aumann's correlated equilibrium permits a mediator to draw a joint action and privately recommend one component to each player \cite{aumann1974subjectivity}.  Its incentive constraints are linear, and existence in finite games admits a duality proof \cite{hart1989existence}; correlated equilibria can also be computed efficiently in broad classes of succinct games \cite{papadimitriou2008computing}.  Symmetric correlated equilibria have been characterized for symmetric $2\times2$ games \cite{amir2017correlated}, while exchangeable equilibrium places an additional conditional-i.i.d. requirement between symmetric Nash and symmetric correlated equilibrium \cite{stein2013exchangeable}.  Tewolde et al. observe that a correlation device can preserve symmetry while coordinating on high-payoff action profiles, but deliberately focus on Nash equilibrium.  We give a general transfer theorem for approximate CE and CCE under arbitrary relabeling groups, prove that welfare-optimal invariant CE exist in every exactly invariant game, and derive a profile-orbit LP.  Thus correlation provides a tractable general-sum endpoint for the same strategic projection used in our Nash analysis.

\paragraph{Symmetry-aware equilibrium solvers and coordination.}
Other-Play exploits known relabeling symmetries to avoid arbitrary conventions in zero-shot coordination \cite{hu2020otherplay}.  Learned normal-form-game solvers have also enforced permutation equivariance when predicting Nash, correlated, and coarse-correlated equilibria \cite{marris2022turbocharging}, and NfgTransformer uses the same representation symmetry for equilibrium solving and deviation-gain prediction \cite{liu2024nfgtransformer}.  These approaches learn a policy or solver across transformed games.  Our goal is complementary: for one explicitly represented game and one supplied group, compute how valid that group is, identify a local certificate when it is not exact, and return an equilibrium whose incentive error in the original game is provably bounded.
